\subsection{Формула полной вероятности. Формула Байеса}
\subsubsection{Введение в теорию групп}
Иногда для нахождения вероятности некоторого события необходимо рассмотреть несколько случаев, которые, в свою очередь, тоже являются событиями. Обозначим такие события как $H_1, H_2, \ldots, H_k$. Обычно их называем гипотезами (отсюда и обозначение через букву $H$). Пусть $H_1 \sqcup H_2 \sqcup \ldots \sqcup H_k = \Omega$.

\begin{definition}[Полная группа событий]
События $H_1, H_2, \ldots, H_k$ образуют полную группу событий, если они попарно несовместны и их объединение -- достоверное событие.
\end{definition}

\subsubsection{Формула полной вероятности}

\begin{theorem}[Формула полной вероятности]
Пусть $H_1, H_2, \ldots, H_k$ -- события, образующие полную группу. Тогда для произвольного события $A$ верно:
\[
\mathbb{P}(A) = \sum_{i=1}^{k} \mathbb{P}(A \mid H_i) \, \mathbb{P}(H_i).
\]
\end{theorem}

\textit{Доказательство:} Рассмотрим набор событий $AH_1, AH_2, \ldots, AH_k$. Очевидно, что если $i \neq j$, то $AH_i \cap AH_j = \varnothing$. Действительно,
\[
AH_i \cap AH_j = A \cap H_i \cap H_j = A \cap (H_i \cap H_j) = A \cap \varnothing = \varnothing.
\]
При этом
\[
AH_1 \sqcup AH_2 \sqcup \ldots \sqcup AH_k = A(H_1 \sqcup H_2 \sqcup \ldots \sqcup H_k) = A\Omega = A.
\]
Тогда из свойств вероятности:
\[
\mathbb{P}(A) = \mathbb{P}(AH_1 \sqcup AH_2 \sqcup \ldots \sqcup AH_k) = \sum_{i=1}^{k} \mathbb{P}(AH_i) = \sum_{i=1}^{k} \mathbb{P}(A \mid H_i) \, \mathbb{P}(H_i). \hfill \blacksquare
\]

\begin{notice}
Заметим, что формула полной вероятности работает и для счётных полных групп, и доказательство будет аналогичным.
\end{notice}

\subsubsection{Формула Байеса}

\begin{theorem}[Формула Байеса]
Пусть $H_1, H_2, \ldots, H_k$ -- события, образующие полную группу. Тогда для произвольного возможного события $A$ и $n \in \overline{1,k}$ верно
\[
\mathbb{P}(H_n \mid A) = \frac{\mathbb{P}(A \mid H_n) \cdot \mathbb{P}(H_n)}{\sum_{i=1}^{k} \mathbb{P}(A \mid H_i) \, \mathbb{P}(H_i)}.
\]
\end{theorem}

\textit{Доказательство:} Формула доказывается с помощью формулы полной вероятности и определения условной вероятности. Действительно:
\[
\mathbb{P}(H_n \mid A) = \frac{\mathbb{P}(AH_n)}{\mathbb{P}(A)} = \frac{\mathbb{P}(A \mid H_n) \cdot \mathbb{P}(H_n)}{\sum_{i=1}^{k} \mathbb{P}(A \mid H_i) \, \mathbb{P}(H_i)}. \hfill \blacksquare
\]


Обычно формулу Байеса интерпретируют так: предположим, что мы знаем вероятности всех исходов $\mathbb{P}(H_i)$, $i \in \overline{1,k}$. Такие вероятности называют \textbf{априорными} (т.е. <<полученными до опыта>>). Нам становится известно, что некоторое событие $A$ произошло, т.е. мы получаем дополнительную информацию. Из-за этого вероятности гипотез <<меняются>> и спрашивается, чему же равны вероятности $\mathbb{P}(H_i \mid A)$, называемыми \textbf{апостериорными} (т.е. полученными <<после опыта>>).


\subsubsection{Цепи Маркова}
\begin{definition}[Цепь Маркова]
Цепь Маркова -- последовательность случайных событий с конечным или счётным числом исходов, где вероятность наступления каждого события зависит только от состояния, достигнутого в предыдущем событии.
\end{definition}

Чтобы описать такую цепь, используют граф переходов между состояниями, где вершины -- состояния, а ориентированные рёбра -- возможные переходы.